\section{\texorpdfstring{\texttt{structure}: bundling data together}{structure: bundling data together}}\label{sec:02-functions-and-structures:01-structure-basics:1}

Chapter 1 fixed what a single term and a single function look like, down to the calculus of constructions itself. But a function of several arguments, as the chapter introduction just noted, is really a chain of one-argument functions --- \texttt{add : Nat → Nat → Nat} never actually receives two numbers at once. That is fine for arguments that stay independent, but plenty of the mathematics this book will build --- a point in the plane, later a group's carrier and operation and axioms together --- needs several pieces of data to travel \emph{together}, as one value, not as separate arguments a caller has to remember to keep in sync. \texttt{structure} is Lean's answer to that need.

A \texttt{structure} groups several pieces of data under one name. We will use this Lean feature constantly once we define groups and rings.

\begin{lstlisting}
structure Point where
  x : Nat
  y : Nat

def origin : Point := { x := 0, y := 0 }

#eval origin.x        -- 0
\end{lstlisting}

Key points:

\begin{itemize}
\tightlist
\item
  \texttt{Point.mk} is the automatically generated \textbf{constructor}. This is the function that actually builds a \texttt{Point} out of an \texttt{x} and a \texttt{y}. \texttt{\{ x := ..., y := ... \}} is shorthand for \texttt{Point.mk ... ...}. It names each field so the order cannot be mixed up.
\item
  There is an even shorter way to write this: \texttt{⟨0, 0⟩} builds the exact same \texttt{Point}. It simply lists the values in field-declaration order instead of naming them. Read \texttt{⟨\_\allowbreak{}, \_\allowbreak{}⟩} as \textbf{``here are the pieces, in order --- the constructor is to be inferred from context.''} Lean can always determine it, because the \emph{expected type} (here, \texttt{Point}, from \texttt{def origin : Point := ...}) indicates exactly which constructor and which fields must be filled in. This is also where the official name comes from: it is called the \textbf{anonymous constructor}, because \texttt{Point.mk} is never written explicitly --- it remains anonymous, and Lean infers it from context. \texttt{⟨\_\allowbreak{}, \_\allowbreak{}⟩} recurs constantly from Chapter 3 onward, for proofs as much as for data.
\item
  \texttt{p.x} is \textbf{field projection} notation, shorthand for \texttt{Point.x p}. \texttt{origin.x} above is exactly this projection applied to \texttt{origin}.
\end{itemize}

\begin{lstlisting}
def shift (p : Point) (dx dy : Nat) : Point :=
  { x := p.x + dx, y := p.y + dy }

#eval (shift origin 3 4).y   -- 4
\end{lstlisting}

\texttt{shift} shows a structure used on \emph{both} sides of a function: it takes a \texttt{Point} in (reading its fields back out with the same \texttt{p.x}/\texttt{p.y} projection notation) and builds a new one via \texttt{\{ x := ..., y := ... \}}, the same field-naming syntax \texttt{origin} used above.

Note also that structures can bundle \emph{proofs} alongside data, not just data. This is exactly how a group will be defined later --- a carrier type, an operation, and proofs that the operation satisfies the group axioms, all in one \texttt{structure}.

\begin{mathreading}

\texttt{structure Point where x : Nat; y : Nat} is the Cartesian product \(\mathrm{Point} = \mathbb{N} \times \mathbb{N}\), with \texttt{x} and \texttt{y} playing the role of the two projections \(\pi_1, \pi_2 : \mathrm{Point} \to \mathbb{N}\). Here \texttt{p.x} computes \(\pi_1(p)\), \texttt{p.y} computes \(\pi_2(p)\), and \texttt{\{ x := ..., y := ... \}} builds an element via the universal property of the product (a pair of maps into the factors determines a unique map into the product). More generally, a \texttt{structure} with fields of types \(A_1, \ldots, A_n\) is the \(n\)-fold product \(A_1 \times \cdots \times A_n\). Once fields are allowed to be \emph{proofs} (propositions), the same construction becomes a \textbf{subset cut out by conditions}: a structure \texttt{\{ data : D, proof : P data \}} is the dependent pair (subset) \(\{\, d \in D \mid P(d) \,\}\), categorically a subobject of \(D\).

\end{mathreading}

Here is the universal property itself, in this box's own notation: the unique \(h\) making both triangles commute is exactly what \texttt{\{ x := ..., y := ... \}} constructs when \(f\) and \(g\) are \(p\)'s two argument expressions.

\begin{center}
% Universal property of the categorical product A x B, in this box's own
% notation: C -> A (f), C -> B (g), and a unique mediating h : C -> A x B
% making both triangles commute via the projections p_A, p_B.
% Source: 02-functions-and-structures/01-structure-basics.md, "Categorical
% product (universal property)" Recall entry / Mathematical reading box.
\begin{tikzcd}
  & C \arrow[dl, "f"'] \arrow[dr, "g"] \arrow[d, dashed, "\exists ! h"] & \\
  A & A \times B \arrow[l, "p_A"] \arrow[r, "p_B"'] & B
\end{tikzcd}

\end{center}

\subsection{Sources, quoted}\label{sec:02-functions-and-structures:01-structure-basics:2}

Formal definitions and citations for this section, gathered here for reference (full entries in the \hyperref[sec:root:bibliography:1]{Bibliography}):

\begin{itemize}
\tightlist
\item
  \textbf{Record / tuple.} ``The simplest of these is pairs, or more generally tuples, of values \ldots{} the generalization from n-ary tuples to labeled records is equally straightforward'' (\cite{Pierce2002}, §11.6--§11.8). Picture it like this: an unlabeled tuple is a box packed in a fixed order --- you just have to remember ``first slot is the name, second is the age.'' A labeled record is the same box with every slot tagged directly: ``Name: \textbf{\emph{, Age: }}.'' Lean's \texttt{structure} is exactly this tagged version.
\item
  \textbf{Categorical product (universal property).} For a category with a product \(A \times B\) (with projections \(p_A, p_B\)) and any object \(C\) with morphisms \(f : C \to A\), \(g : C \to B\): ``there is exactly one morphism \(h : C \to A \times B\) such that \(f = p_A h\) and \(g = p_B h\)'' (\cite{Pareigis1970}, §1.11, p.~30).
\item
  Lean 4 documentation (\cite{LeanDocs}) --- the constructor/projection/anonymous-constructor mechanics described above.
\end{itemize}
